\section{Homework 01: Panoramas, Features, and Calibration}

This homework investigated image alignment tasks on paired RGB and thermal
drone images. The available scenes were \texttt{ellipse}, \texttt{FH3},
\texttt{forest}, and \texttt{hut}. They differ strongly in camera motion,
scene structure, and overlap, which made them suitable for comparing panorama
stitching, cross-modal feature matching, intrinsic calibration, and relative
pose estimation. The RGB images have a larger field of view, while the thermal
camera shows stronger distortion and less stable local texture. Therefore, the
evaluation is based not only on numerical success values, but also on visual
inspection and geometric consistency checks.

\subsection{Panorama Stitching}

The panorama experiment was carried out on the \texttt{ellipse} scene. This
scene is the most appropriate sequence for panorama construction because the
drone mostly rotated around the vertical axis while staying close to one
position. This motion fits the panorama assumption better than the other
scenes, where translation, depth variation, or vegetation produce stronger
parallax. Seven RGB frames and seven thermal frames were used.

The final implementation did not use OpenCV's automatic \texttt{Stitcher}.
Instead, all frames were stitched manually in a controlled pipeline. First, the
images were resized to a common processing scale. Thermal frames were converted
to grayscale and contrast-enhanced with CLAHE, because raw thermal intensities
provide less stable local texture than RGB images. Second, each frame was
projected onto a cylinder. This step is important because a yaw rotation of the
camera becomes approximately a horizontal translation after cylindrical
projection. A simple affine or homography chain on the original images produced
more visible ghosting and was therefore less suitable for this scene.

After projection, neighboring frames were aligned with SIFT features and a
small RANSAC procedure over 2D translation offsets. The RGB alignment was very
stable: all six neighboring pairs were accepted, with 920 to 1333 inliers and
inlier ratios between \(0.899\) and \(0.950\). The estimated horizontal
translations ranged from approximately \(172\) to \(255\) pixels, while the
vertical offsets stayed small. The thermal alignment was also accepted for all
six neighboring pairs, but it was clearly weaker. It produced 246 to 587
inliers and inlier ratios between \(0.344\) and \(0.544\). This difference
matches the visual properties of the data: the RGB frames contain sharper
building texture, while the thermal frames contain fewer repeatable features
and stronger lens distortion.

The pairwise translations were chained into the middle frame as reference. All
cylindrical images were then warped onto one common canvas and blended with
distance-transform feather weights. The resulting RGB panorama has a size of
\(582 \times 2143\) pixels and a non-black image fraction of \(94.01\,\%\). It
is visually coherent, and the building structure remains mostly consistent.
Small artifacts can still appear because the implementation does not include
bundle adjustment, advanced seam finding, exposure compensation, or multi-band
blending.

The thermal panorama has a size of \(841 \times 2884\) pixels and a non-black
image fraction of \(94.20\,\%\). It is usable and visually coherent after CLAHE,
cylindrical projection, translation-based stitching, and feather blending.
However, the result must be interpreted more carefully than the RGB panorama.
The focal value for the cylindrical projection is only approximate, the thermal
camera is more distorted, and the image content contains fewer stable local
details. Therefore, the thermal panorama demonstrates that the method can be
applied to the thermal sequence, but it is less geometrically reliable than the
RGB result.

\subsection{Thermal-to-RGB Image Features}

For thermal-to-RGB matching, all four scenes were evaluated with three feature
setups: SIFT on CLAHE-enhanced grayscale images, ORB on normalized grayscale
images, and AKAZE on edge-enhanced images. The matching pipeline used
descriptor matching with Lowe's ratio test, mutual matching, unique RGB target
points, and a RANSAC homography as geometric verification. The quality was
measured with the number of detected keypoints, remaining descriptor matches,
RANSAC inliers, inlier ratio, median reprojection error, spatial spread of
matches, and visual inspection of the resulting match lines.

The experiment shows that matching thermal to RGB images is much harder than
matching two RGB images. RGB texture is often not visible in the thermal image,
and thermal intensity changes do not necessarily correspond to visible image
edges. High keypoint counts therefore do not directly translate into many
reliable matches. SIFT and ORB often detected thousands of keypoints, but only a
small number survived mutual matching and RANSAC. AKAZE on edge-enhanced images
was the most useful descriptor setup overall, because structural boundaries are
more likely to appear in both modalities.

The best result was obtained on the \texttt{ellipse} scene with AKAZE, where
the strongest pair produced 33 descriptor matches and 18 RANSAC inliers. This
was classified as the only clearly good case in the experiment, although the
visual result was still fragile. The \texttt{FH3} and \texttt{hut} scenes
produced weak but usable examples, with best AKAZE results of 6 and 7 inliers,
respectively. The \texttt{forest} scene did not work reliably. Its repetitive
vegetation texture produced too few visually meaningful correspondences, even
when RANSAC found a small consistent subset.

The main conclusion is that thermal-to-RGB feature matching is possible for
structured scenes, but it is not robust enough to be treated as a fully
automatic registration method in this dataset. Descriptor distance alone is not
a sufficient quality measure. A useful evaluation must combine match counts,
RANSAC inliers, reprojection error, inlier spread, and visual inspection.

\subsection{Intrinsic Camera Calibration}

The intrinsic calibration task was treated separately from the provided
ground-truth intrinsics. The goal was to check whether the available HW01 RGB
and thermal images themselves are sufficient to estimate camera matrices and
distortion coefficients. Since the RGB and thermal images come from two
different physical cameras, they would have to be calibrated independently. The
dataset contains 22 RGB images of size \(3840 \times 2160\) pixels and 22
thermal images of size \(1280 \times 1024\) pixels. This number of images would
be reasonable for a calibration experiment only if they showed a known
calibration target from different positions, distances, and viewing angles.

The available images are natural drone scenes rather than calibration-target
images. This is the central limitation. A standard OpenCV calibration workflow
requires known 3D object points, such as checkerboard or circle-grid corner
coordinates in metric units, and corresponding detected 2D image points. The
HW01 scenes contain buildings, vegetation, and other scene structures, but the
real 3D positions, distances, and physical dimensions of these structures are
unknown. Manually clicking image points would therefore not solve the problem,
because the matching 3D object points would still be missing.

As a sanity check, common checkerboard corner layouts were searched in all HW01
images. No valid calibration-pattern candidates were found: the result was
\(0\) candidates for all 22 RGB images and \(0\) candidates for all 22 thermal
images. Consequently, the lists of object points and image points remained
empty, and \texttt{cv2.calibrateCamera} was intentionally not executed. Running
the function without valid 3D-to-2D correspondences would not produce a
meaningful calibration result.

The conclusion is that a real intrinsic calibration is not feasible from the
HW01 images alone. The limitation is not the calibration algorithm, but the
input data. Without known metric calibration geometry, the estimation problem is
underconstrained: many different focal lengths, principal points, and
distortion parameters could explain natural-scene image measurements. Feature
matches between frames are useful for homography estimation, panorama
construction, and relative pose estimation, but they are only 2D-to-2D
correspondences and cannot replace calibration-target measurements.

With valid calibration data, calibration quality would be evaluated with the
RMS reprojection error, per-image residuals, residual plots over the image,
visual undistortion checks on straight lines, and preferably a validation set of
held-out calibration images. In this experiment, no such quality score can be
reported because no calibration target and no known 3D-to-2D correspondences
are available. Therefore, the provided ground-truth intrinsics may be used for
later undistortion or extrinsic estimation, but they should not be presented as
intrinsics estimated from the HW01 image set.

\subsection{Extrinsic Estimation}

The extrinsic notebook estimated relative motion frame by frame. The provided
ground-truth intrinsic matrices and distortion coefficients were used for both
the RGB and thermal cameras, but no ground-truth extrinsic parameters were used.
For each consecutive frame pair, the pipeline detected SIFT features on
CLAHE-enhanced grayscale images, matched descriptors with a mutual Lowe ratio
test, undistorted the matched image points, and estimated an essential matrix
with RANSAC. The relative rotation and translation direction were then recovered
with \texttt{recoverPose}.

The interpretation of these results is limited by monocular geometry. The
translation can only be recovered up to an unknown scale, so the estimated
translation vector is only a direction. In addition, pure rotation, low
parallax, planar structure, and repetitive texture can make the translation
direction unstable even when the number of inliers is high. For this reason,
the evaluation used pose inlier count, inlier ratio, median epipolar error,
median parallax angle, a homography-versus-Essential-Matrix degeneracy check,
and visual inspection.

The \texttt{FH3} sequence was the most consistently usable scene. All four RGB
pairs and all four thermal pairs were classified as usable, with RGB inlier
ratios between \(0.73\) and \(0.91\), and thermal inlier ratios between
\(0.49\) and \(0.67\). The \texttt{forest} sequence contained some poor pairs
with too few matches, but later frame pairs showed strong pose support, for
example 916 RGB pose inliers on the best pair. The \texttt{hut} sequence was
mixed, with some weak pairs and one clearly usable final pair for both cameras.

The \texttt{ellipse} sequence highlights the main ambiguity of the task. Since
the drone mostly rotated around the vertical axis, the rotation can be recovered
more reliably than the translation. The first RGB pair had 1240 pose inliers
and an inlier ratio of \(0.84\), but it was still classified as rotation usable
with weak translation because the motion provides limited parallax and is close
to a panorama setup. Most later \texttt{ellipse} pairs were poor for relative
pose estimation.

Overall, extrinsic estimation is feasible in a weak relative sense when stable
matches and sufficient parallax are available. It should not be interpreted as
a metric drone trajectory, because no absolute scale or ground-truth pose is
available. Rotation estimates are generally more trustworthy than translation
directions, and the final quality judgment has to combine geometric metrics
with visual inspection of the matched inliers.
